\graphicspath{{abstracts/byw/01-briefedition-mueller/img/}{img/}}

\begin{beitrag}{Projekt \glqq Briefedition Wilhelm Müller\grqq{}: WissKI als Editionswerkzeug}{%
Wolfgang Wagner\,\orcidlink{0000-0002-3456-7891} \\
\small \href{mailto:wolfgang.wagner@example.org}{wolfgang.wagner@example.org} \\
\small Universität Stuttgart
}

\section{Projektsteckbrief}
Das im Format \emph{Bring your WissKI} vorgestellte Projekt erschließt den Korrespondenzbestand des Dichters und Bibliothekars Wilhelm Müller (1794--1827) in einer Online-Edition, die zugleich als kommentierte Lesefassung und als angereichertes Forschungsdatensystem nutzbar sein soll. Die Edition umfasst rund 720 überlieferte Briefe von und an Müller; ergänzend werden eine Personen-, eine Orts- und eine Werkdatenbank aufgebaut, die das Korrespondenznetz situieren. Das Projekt läuft seit gut zwei Jahren und befindet sich aktuell in der Mitte seiner Förderlaufzeit; die kritischen Designentscheidungen sind weitgehend getroffen, die wesentliche Erschließungsarbeit ist in vollem Gange.

\section{Aktueller Stand der Arbeit}
Bisher wurden etwa 280 der überlieferten Briefe vollständig transkribiert, mit XML-TEI-Markup versehen und in der WissKI-Edition publiziert. Die zugehörigen Personeneinträge sind weitgehend mit GND-Identifikatoren angereichert, ebenso die Ortseinträge mit GeoNames- und Wikidata-IDs. Die Modellierung folgt einem an CIDOC CRM angelehnten Profil, ergänzt um ein TEI-Subset für die Volltexte. Ein detaillierter Erfahrungsbericht aus den ersten achtzehn Monaten Projektarbeit ist im vergangenen Sommer erschienen \cite{wagner2024-mueller-bericht}.

Die Editionsoberfläche kombiniert eine diplomatische und eine lesefreundliche Sicht auf den Brieftext mit einer kontextuellen Anzeige, die Personen, Orte, Werke und Ereignisse, die im Brief erwähnt werden, durch ein klickbares Overlay zugänglich macht. Diese Anreicherung ist nicht statisch hinterlegt, sondern entsteht zur Laufzeit aus den im Triplestore gespeicherten Beziehungen.

\section{Worüber wir diskutieren möchten}
Im Format \emph{Bring your WissKI} bringen wir drei konkrete Fragestellungen zur Diskussion:

Erstens die Frage der \emph{Granularität der Annotation}: Wie weit soll die Erschließung gehen? Bisher annotieren wir Personen, Orte und Werke; die Aufnahme von Konzepten (\glqq Klassizismus\grqq{}, \glqq Wandermotiv\grqq{}, \glqq Volkslied\grqq{}) ist bisher nicht systematisch erfolgt. Die fachliche Diskussion im Projektteam ist hier nicht abgeschlossen. Wir würden gerne hören, welche Erfahrungen andere Editionsprojekte mit konzeptueller Annotation gemacht haben und ab welchem Detailgrad der Aufwand in einem produktiven Verhältnis zum Nutzen steht.

Zweitens die Frage des \emph{Übergangs vom Forschungs- in den Daueretat}: Das Projekt ist als Drittmittelvorhaben angelegt; nach Ende der Förderung soll die Edition aber langfristig verfügbar bleiben. Wir würden gerne erfahren, welche Strategien sich in vergleichbaren Projekten bewährt haben, um diesen Übergang nicht erst am Projektende zu organisieren.

Drittens die Frage der \emph{Anbindung an externe Editionsplattformen}: Soll die Müller-Edition in eine größere Plattform (TextGrid, CLARIAH) integriert werden, oder ist ein eigenständiger Auftritt vorzuziehen? Wir diskutieren intern unterschiedliche Positionen und würden gerne von den Erfahrungen anderer Editionsprojekte lernen \cite{wagner2025-plattformen}.

\section{Vorgesehenes Format}
Wir möchten unsere Diskussionsbeiträge in einer Mischung aus Live-Demonstration der Editionsoberfläche und offener Gesprächsrunde einbringen. Wir haben eine Auswahl konkreter Editionsstellen vorbereitet, an denen sich die genannten Diskussionsfragen exemplarisch festmachen lassen; wir werden diese Stellen im Plenum zeigen und freuen uns auf die anschließende Diskussion mit den Teilnehmer*innen des Anwender*innentreffens.

\section{Was wir mitbringen, was wir suchen}
Wir bringen mit: konkrete Erfahrungen aus zwei Jahren TEI-XML-WissKI-Integration, ein dokumentiertes Datenmodell, eine produktiv laufende Edition mit etwa 280 erschlossenen Briefen, und eine Reihe technischer Lösungen, die wir gerne zur Diskussion stellen. Wir suchen: einen Austausch mit Editionsprojekten, die ähnliche Fragen bearbeiten, mit besonderem Interesse an Projekten, die sich mit konzeptueller Annotation und mit der Frage der langfristigen Verstetigung beschäftigen.

\end{beitrag}
